\sectionkicker{Source-backed technical notes}
\subsection*{モデルを動かす技術が競争軸になった――QLoRA、FlashAttention-2、vLLM: Technical Notes}
本欄は記事本文で圧縮した一次資料上の情報を、比較・再検証しやすい形で整理したものである。時系列、確認済みの主張、留意点、一次資料URLを掲載する。
\smallskip
\noindent{\small\textit{Technical Notesでは、一次資料から確認した公開・提供・研究上の事実を記載する。数値・能力評価や提供元・著者による比較表現は、独立再現済みの結果ではなく帰属claimとして扱う。}}\par
\medskip
\subsection*{Theme at a glance}
\addcontentsline{toc}{subsection}{Theme at a glance}
\begin{center}
\footnotesize
\begin{tabularx}{\linewidth}{@{}p{0.20\linewidth}p{0.10\linewidth}p{0.14\linewidth}X@{}}
\toprule
資料 & 位置づけ & 種別 & 時系列 \\
\midrule
vLLM / PagedAttention & 主要資料 & Framework & 2023-06-20, 2023-08-02, 2023-09-12 \\
QLoRA & 主要資料 & 論文 & 2023-05-23 \\
FlashAttention-2 & 主要資料 & 論文 & 2023-07-17 \\
Phi-2 & 補足資料 & モデル & 2023-12-12 \\
\bottomrule
\end{tabularx}
\end{center}
\normalsize

\begin{technicalnote}{vLLM / PagedAttention}{主要資料}
\begingroup
% reader-facing Technical Notes break policy
\widowpenalty=10000
\clubpenalty=10000
\displaywidowpenalty=10000
\begin{tabularx}{\linewidth}{@{}>{\bfseries}p{0.22\linewidth}X@{}}
組織 & vLLM Project \\
種別 & Framework \\
時系列 & 2023-06-20, 2023-08-02, 2023-09-12 \\
\end{tabularx}
\smallskip
{\bfseries 一次資料から整理したtechnical points}
\begin{itemize}[leftmargin=1.5em,itemsep=0.35em]
\item \textbf{一次情報で確認できる事実}: vLLMの論文では、PagedAttentionはOSのvirtual memory/pagingに着想を得てKV cacheをblock単位で管理し、vLLMはKV-cache memory wasteをほぼゼロにし、request内外でKV cacheを共有できるよう設計されていると著者が報告する。
\end{itemize}
\begin{minipage}{\linewidth}
% reader-facing Technical Notes generic boundary/source tail group
\begin{samepage}
{\bfseries 一次資料}
\begingroup\sloppy
\Urlmuskip=0mu plus 2mu\relax
\begin{itemize}[leftmargin=1.5em,itemsep=0.25em]
\item {\scriptsize\url{http://arxiv.org/abs/2309.06180v1}}
\item {\scriptsize\url{https://github.com/vllm-project/vllm/releases/tag/v0.1.3}}
\item {\scriptsize\url{https://vllm.ai/blog/2023-06-20-vllm}}
\end{itemize}
\endgroup
\end{samepage}
\end{minipage}
\endgroup
\end{technicalnote}
\medskip

\begin{technicalnote}{QLoRA}{主要資料}
\begingroup
% reader-facing Technical Notes break policy
\widowpenalty=10000
\clubpenalty=10000
\displaywidowpenalty=10000
\begin{tabularx}{\linewidth}{@{}>{\bfseries}p{0.22\linewidth}X@{}}
組織 & authors \\
種別 & 論文 \\
時系列 & 2023-05-23 \\
\end{tabularx}
\smallskip
{\bfseries 一次資料から整理したtechnical points}
\begin{itemize}[leftmargin=1.5em,itemsep=0.35em]
\item \textbf{一次情報で確認できる事実}: QLoRAはfrozen 4-bit quantized pretrained modelを通してgradientをbackpropagateし、学習対象をlow-rank adapter weightsへ限定する。論文はnormally distributed weights向けの4-bit NormalFloat (NF4)と、NVIDIA unified memoryを介してoptimizer-state pressureを処理するpaged optimizersをmemory-saving mechanismsとして導入する。
\end{itemize}
\begin{minipage}{\linewidth}
% reader-facing Technical Notes generic boundary/source tail group
\begin{samepage}
{\bfseries 一次資料}
\begingroup\sloppy
\Urlmuskip=0mu plus 2mu\relax
\begin{itemize}[leftmargin=1.5em,itemsep=0.25em]
\item {\scriptsize\url{http://arxiv.org/abs/2305.14314v1}}
\end{itemize}
\endgroup
\end{samepage}
\end{minipage}
\endgroup
\end{technicalnote}
\medskip

\begin{technicalnote}{FlashAttention-2}{主要資料}
\begingroup
% reader-facing Technical Notes break policy
\widowpenalty=10000
\clubpenalty=10000
\displaywidowpenalty=10000
\begin{tabularx}{\linewidth}{@{}>{\bfseries}p{0.22\linewidth}X@{}}
組織 & authors \\
種別 & 論文 \\
時系列 & 2023-07-17 \\
\end{tabularx}
\smallskip
{\bfseries 一次資料から整理したtechnical points}
\begin{itemize}[leftmargin=1.5em,itemsep=0.35em]
\item \textbf{一次情報で確認できる事実}: authorsの選定済み一次資料では、FlashAttention-2 は exact attention の IO-aware algorithm を改良し、non-matmul FLOPs を減らすとともに、sequence-length dimension も含めて thread blocks 間の work partitioning を改善する。 単一 thread block 内では warps 間の work distribution を変更して shared-memory communication を減らし、GPU occupancy と実効 throughput の改善を狙う。 論文中の speedup / FLOPs utilization は著者実験の結果として扱い、特定GPU・sequence条件を超えて一般化しない。
\end{itemize}
\begin{minipage}{\linewidth}
% reader-facing Technical Notes generic boundary/source tail group
\begin{samepage}
{\bfseries 一次資料}
\begingroup\sloppy
\Urlmuskip=0mu plus 2mu\relax
\begin{itemize}[leftmargin=1.5em,itemsep=0.25em]
\item {\scriptsize\url{http://arxiv.org/abs/2307.08691v1}}
\end{itemize}
\endgroup
\end{samepage}
\end{minipage}
\endgroup
\end{technicalnote}
\medskip

\begin{technicalnote}{Phi-2}{補足資料}
\begingroup
% reader-facing Technical Notes break policy
\widowpenalty=10000
\clubpenalty=10000
\displaywidowpenalty=10000
\begin{tabularx}{\linewidth}{@{}>{\bfseries}p{0.22\linewidth}X@{}}
組織 & Microsoft Research \\
種別 & モデル \\
時系列 & 2023-12-12 \\
\end{tabularx}
\smallskip
{\bfseries 一次資料から整理したtechnical points}
\begin{itemize}[leftmargin=1.5em,itemsep=0.35em]
\item \textbf{一次情報で確認できる事実}: Microsoft Researchの選定済み一次資料では、Microsoft Research は Phi-2 を 2.7B-parameter language model として公開し、training data quality を重視した small-model 系列の継続として位置づけた。 公開説明では synthetic datasets と carefully selected web data を含む textbook-quality data を学習設計の中心に置き、model scale だけでなく data curation を性能要因として強調している。 benchmark 上の大型modelとの比較は Microsoft Research の評価として扱い、独立再現済みの優位性とはみなさない。
\end{itemize}
\begin{minipage}{\linewidth}
% reader-facing Technical Notes generic boundary/source tail group
\begin{samepage}
{\bfseries 一次資料}
\begingroup\sloppy
\Urlmuskip=0mu plus 2mu\relax
\begin{itemize}[leftmargin=1.5em,itemsep=0.25em]
\item {\scriptsize\url{https://www.microsoft.com/en-us/research/blog/phi-2-the-surprising-power-of-small-language-models/}}
\end{itemize}
\endgroup
\end{samepage}
\end{minipage}
\endgroup
\end{technicalnote}
\medskip
